\section{Canonical endpoint normalization}\label{tsc:normalization}
Write $\TT=\RR/\ZZ$ and $e(x)=\exp(2\pi\ii x)$. Fourier coefficients use
\[
 \widehat f(n)=\langle f,e(-nx)\rangle.
\]
For $n\ne0$ fix
\begin{equation}\label{tsc:branch}
 (2\pi\ii n)^{-u}=(2\pi|n|)^{-u}
       \exp\!\left(-\frac{\pi\ii u}{2}\sgn n\right).
\end{equation}
All shifts below are real modulo one. Values assigned at isolated seam
points do not affect ordinary integrals. Spectral derivatives are denoted
$\zeta^{[j]}(s,x)=\partial_s^j\zeta(s,x)$ and
$\Li_s^{[j]}(z)=\partial_s^j\Li_s(z)$. Parenthesized superscripts on
$\gamma_m$ or $\psi$ denote derivatives in the spatial argument.

We deliberately use the spectral coordinate $u=1-s$:
\begin{equation}\label{tsc:stieltjes-convention}
 \zeta(1-u,x)=-\frac1u+
     \sum_{m\ge0}\frac{\gamma_m(x)}{m!}u^m,
 \qquad \gamma_0(x)=-\psi(x).
\end{equation}
This changes no Stieltjes convention: it is the usual Laurent expansion
at $s=1$ after substituting $s=1-u$. The shift relation gives
\begin{equation}\label{tsc:gamma-local}
 \gamma_m(x)=\frac{\log^m x}{x}+\gamma_m(1+x),\qquad x>0.
\end{equation}

\begin{definition}[Scale-one periodic Stieltjes extension]\label{tsc:Cdef}
For a smooth periodic test function $\varphi$, define
\begin{equation}\label{tsc:Cfp}
 \langle C_m,\varphi\rangle=
 \lim_{\varepsilon\downarrow0}\left\{
   \int_{\varepsilon}^{1}\gamma_m(x)\varphi(x)\dd x
    +\frac{\varphi(0)}{m+1}\log^{m+1}\varepsilon\right\}.
\end{equation}
The coordinate at the singularity is exactly $x$, without a rescaling or
an additional subtraction length. Set $C_{m,a}(x)=C_m(x-a)$.
\end{definition}
Existence follows from \eqref{tsc:gamma-local}: after subtracting
$\varphi(0)\log^m x/x$, the integrand is $O(1+|\log x|^m)$.

\begin{lemma}[Entire normalized Hurwitz family]\label{tsc:entire-B}
The Fourier prescription
\begin{equation}\label{tsc:Bdef}
 \calB_u(x)=\sum_{n\ne0}(2\pi\ii n)^{-u}e(nx)
\end{equation}
defines an entire distribution-valued function of $u\in\CC$, with mean
zero and singular support contained in $0$. On $0<x<1$,
\begin{equation}\label{tsc:Bhurwitz}
 \calB_u(x)=\frac{\zeta(1-u,x)}{\Gamma(u)},
\end{equation}
with removable values interpreted analytically. Near $u=0$,
\begin{equation}\label{tsc:Zlaurent}
 \Gamma(u)\calB_u=\frac{\delta_0-1}{u}
          +\sum_{m\ge0}\frac{C_m}{m!}u^m.
\end{equation}
In particular, $\calB_0=\delta_0-1$ and $\langle C_m,1\rangle=0$.
\end{lemma}
\begin{proof}
On a compact set of spectral parameters, the coefficients in
\eqref{tsc:Bdef}, and every spectral derivative of them, grow at most
polynomially in $|n|$, with additional powers of $\log|n|$. The Fourier
coefficients of a smooth test function decrease faster than every power.
Their pairing is therefore normally convergent, including all spectral
derivatives. This proves entire dependence in distributions.

For $\Re u>1$, the classical Hurwitz Fourier expansion
\cite[Eq.~25.11.9]{tsc:dlmf-zeta} gives \eqref{tsc:Bhurwitz}; both sides
continue analytically away from the seam. The right side is smooth there
and entire in $u$, since the pole of $\zeta(1-u,x)$ at zero is canceled by
$1/\Gamma(u)$. This also proves the assertion about singular support.

For the Laurent normalization, split a test-function integral near zero
and use
\[
 \zeta(1-u,x)=x^{u-1}+\zeta(1-u,1+x).
\]
The first term has residue $\varphi(0)$ at $u=0$; the second has residue
$-\int\varphi$. The regular coefficients of the first term are precisely
the scale-one finite parts of $\log^m x/x$. The regular coefficients of
the second are smooth. This proves \eqref{tsc:Zlaurent} and the stated
normalization. Mean zero follows either from the Fourier prescription or
by comparing Laurent coefficients of its zero Fourier coefficient.
\end{proof}

\subsection{Products at separated singularities}
Let $a_1,a_2,a_3$ be pairwise distinct in $\TT$. Near $a_i$, the other two
translated factors are smooth. A partition of unity therefore defines a
unique product of these distributions: multiply the singular factor by the
smooth factors on each such neighborhood and use ordinary multiplication
elsewhere. The definition is independent of the partition. It also gives
joint entire dependence for products of the $\calB_{u_i}$, because locally
only one distribution is multiplied by smooth parameter-dependent factors.
No assertion about products at a common singularity is needed.

\begin{definition}[Separated triple finite part]\label{tsc:Jdef}
For $\mathbf m=(m_1,m_2,m_3)\in\ZZ_{\ge0}^3$, put
\begin{equation}\label{tsc:J}
 \calJ_{\mathbf m}(\mathbf a)=
  \left\langle\prod_{i=1}^3 C_{m_i,a_i},1\right\rangle
  =\FP\int_0^1\prod_{i=1}^3
       \gamma_{m_i}(\{x-a_i\})\dd x.
\end{equation}
At each seam $a_i$, omit the right-hand arc $0<t<\varepsilon_i$ and add
\begin{equation}\label{tsc:triple-counterterm}
 \frac{\log^{m_i+1}\varepsilon_i}{m_i+1}
       \prod_{j\ne i}\gamma_{m_j}(\{a_i-a_j\}).
\end{equation}
The three cutoffs may tend to zero independently.
\end{definition}
The product of the two smooth factors differs from its seam value by
$O(t)$, so the same local subtraction proof establishes existence. The
one-sided nature of these counterterms must be distinguished from the
symmetric principal values used later for cotangent poles.
